\makeabstract
{ % #1: Title
Examining the effect of S-palmitoylation on Na\textsubscript{V}1.2 localization and its interaction with ankyrins
}
{ % #2: Authorship
Madeline Gold\textsuperscript{1},
Kalynn Harvey\textsuperscript{2},
Paul Jenkins\textsuperscript{2,3}
}
{ % #3: Affiliations (use \\ to start a newline)
\textsuperscript{1} Amherst College, Amherst, MA\\
\textsuperscript{2} Department of Pharmacology, University of Michigan Medical School, Ann Arbor, MI\\
\textsuperscript{3} Department of Psychiatry, University of Michigan Medical School, Ann Arbor, MI
}
{ % #4: Purpose
Voltage-gated sodium channels (Na\textsubscript{V}s) are critical mediators of action potential generation and propagation in neurons. Improper localization of these channels impairs their function, and neuronal excitability, as a result. Na\textsubscript{V}s are anchored to specific excitable regions of the neuronal membrane by members of the ankyrin scaffolding protein family. Mutations of the genes encoding ankyrins and certain Na\textsubscript{V}s that disrupt localization have been implicated in a number of neurodevelopmental conditions. In particular, loss-of-function mutations of the \textit{SCN2A} gene, which encodes the sodium channel Na\textsubscript{V}1.2, represent one of the most prominent single-gene risk factors for autism spectrum disorder (ASD). The exact mechanism of interaction and localization between Na\textsubscript{V}1.2 and its ankyrin binding partners has not yet been confirmed, but studies of other Na\textsubscript{V}1 family members have identified the post-translational modification of S-palmitoylation as a potential regulator of Na\textsubscript{V}/ankyrin binding. It has been shown that palmitoylation-deficient Na\textsubscript{V}1.2 has altered localization early in development in cultured mouse neurons. It has also been shown that Na\textsubscript{V}1.6 is palmitoylated at three distinct sites, two of which are located in the cytoplasmic loop (II-III loop) between the second and third homologous domains of the channel. This loop contains the ankyrin-binding domain, making it of particular importance.
}
{ % #5: Methods
Due to the conservation of cysteine residues across the Na\textsubscript{V}1 family, we sought to create a mutant of the II-III loop of Na\textsubscript{V}1.2 in which the cysteine sites were replaced with alanine residues to examine the effects of palmitoylation on Na\textsubscript{V}1.2. Palmitoylation status of the wild type (WT) versus mutant II-III loop was assessed via the acyl biotin exchange (ABE) assay. Impacts of palmitoylation on the Na\textsubscript{V}1.2/ankyrin complex were assessed via co-immunoprecipitation.
}
{ % #6: Results
We confirmed production of the mutant II-III loop via DNA gel and sequencing. We saw that the mutant II-III loop does not appear to undergo S-palmitoylation. We also saw interaction between the II-III loop and ankyrins B and G, although this occurred with both the WT and mutant loop.
}
{ % #7: Conclusion
Future studies will be conducted to optimize ABE and immunoprecipitation protocols so as to better assess palmitoylation-dependent interactions between the Na\textsubscript{V}1.2 II-III loop and ankyrin binding partners and to explore the functional impacts of palmitoylation-dead Na\textsubscript{V}1.2 in neurons.
}
 \newpage